% ===========================================================
%  CHAPTER IV — THEORETICAL FRAMEWORK OF DRL APPLIED TO MUSCLE CONTROL
% ===========================================================
\chapter{Theoretical Framework of DRL Applied to Muscle Control}
\label{chap:drl}

% -----------------------------------------------------------
\section{MDP Formulation}

The state $s_t \in \mathcal{S} \subset \R^{38}$ is the observation
vector~\eqref{eq:obs}. The action $a_t \in \mathcal{A} = [0,1]^9$ is the vector of
muscle excitations. The transition function $P(s_{t+1}|s_t,a_t)$ is deterministic
in the standard configuration (physics simulation) and stochastic under Domain
Randomization.

The Bellman equation for the action-value function $Q^\pi$ is:
\begin{equation}
  Q^\pi(s,a) = \E\!\left[r(s,a) + \gamma\,Q^\pi(s', \pi(s'))\right]
  \label{eq:bellman}
\end{equation}
The dimensionality of the problem is notable: the state space $\R^{38}$ combined
with the action space $[0,1]^9$ and the nonlinear dynamics of the Hill model make
this a demanding benchmark for \DRL{} algorithms.

\begin{proposition}[Absence of strict Markov guarantee]
Strictly speaking, the observation $o_t$ described by~\eqref{eq:obs} is Markovian
only under the following assumptions: (i)~the second derivatives of position
(accelerations) are sufficiently reconstructible from velocities; (ii)~the
electromechanical delay in activation dynamics (Equation~\eqref{eq:activation}) is
a deterministic function of the observed variables. The \textit{frame-stacking}
paradigm (stacking $k$ consecutive observations) was not required here, consistent
with standard practice~\parencite{Caggiano2022}.
\end{proposition}

% -----------------------------------------------------------
\section{Neural Network Architectures}

All algorithms use MLPs with the same base architecture
(Table~\ref{tab:nn_arch}), ensuring a fair comparison:

\begin{table}[H]
  \centering
  \caption{Network sizes, counted from the saved models. The actor input is
           $\R^{38}$; the critic input is $\R^{38+9} = \R^{47}$.}
  \label{tab:nn_arch}
  \rowcolors{2}{lightblue!30}{white}
  \small
  \setlength{\tabcolsep}{5pt}
  \begin{tabularx}{\textwidth}{@{} L{3.0cm} X C{2.0cm} C{2.0cm} @{}}
    \toprule
    \rowcolor{lightblue}
    \textbf{Component} & \textbf{Architecture} & \textbf{Activ.} & \textbf{Params.} \\
    \midrule
    Actor, \SAC        & $[38 \to 256 \to 256 \to 2{\times}9]$ & ReLU + tanh & 80\,402 \\
    Actor, \TDIII/\DDPG & $[38 \to 256 \to 256 \to 9]$      & ReLU + tanh & 78\,089 \\
    Critic (one)       & $[47 \to 256 \to 256 \to 1]$        & ReLU        & 78\,337 \\
    Twin critics, \SAC/\TDIII & two of the above              & ReLU        & 156\,674 \\
    \PPO{} extractor   & separate policy and value trunks      & Tanh        & 151\,552 \\
    \PPO{} heads       & $256 \to 9$ and $256 \to 1$          & ---         & 2\,570 \\
    \midrule
    \multicolumn{3}{@{}l}{\textbf{Complete saved object} --- \SAC}    & \textbf{393\,750} \\
    \multicolumn{3}{@{}l}{\hphantom{\textbf{Complete saved object}} --- \TDIII} & 469\,526 \\
    \multicolumn{3}{@{}l}{\hphantom{\textbf{Complete saved object}} --- \DDPG} & 312\,852 \\
    \multicolumn{3}{@{}l}{\hphantom{\textbf{Complete saved object}} --- \PPO}  & 154\,131 \\
    \bottomrule
  \end{tabularx}
  \normalsize
\end{table}

Three points follow from \cref{tab:nn_arch}. The \SAC{} actor is larger than the
\TDIII/\DDPG{} actor because it emits a mean \emph{and} a log-standard-deviation
per muscle, $2\times9$ outputs rather than $9$. \SAC{} and \TDIII{} carry twin
critics and their targets, which is where most of their parameter count sits.
\PPO{} builds separate policy and value trunks rather than sharing one, so its
extractor is two $[38 \to 256 \to 256]$ stacks.

Only the actor is needed at deployment: \SI{80402}{} parameters for \SAC{} and
\SI{78089}{} for the others, against the \SI{393750}{} of the complete saved
object. The \SI{1.5}{\mega\byte} footprint quoted in
\cref{sec:compute_efficiency} is the full object in single precision, which is
the conservative figure.

An earlier version of this table gave 76\,809 / 76\,809 / 79\,617 / 76\,545.
None matched, and the first two could not both be right in any case, since the
two actors differ in output dimension.

Weights are initialised using the orthogonal method for PPO~\parencite{Schulman2017}
and uniform Xavier initialisation~\parencite{Glorot2010} for off-policy algorithms.
Observation normalisation is handled by an exponential moving average layer
(EMA, $\alpha = 0.001$) via Stable-Baselines3's \texttt{VecNormalize}.

% -----------------------------------------------------------
\section{Theoretical Comparison of Algorithms}

\begin{table}[H]
  \centering
  \caption{Theoretical comparison of the implemented \DRL{} algorithms.}
  \label{tab:algo_compare}
  \rowcolors{2}{lightblue!30}{white}
  \small
  \begin{tabularx}{\textwidth}{@{} L{1.6cm} L{1.6cm} L{2.2cm} L{2.5cm} X L{2.2cm} @{}}
    \toprule
    \rowcolor{lightblue}
    \textbf{Algo.} & \textbf{Type} & \textbf{Policy} & \textbf{Exploration} &
    \textbf{Advantages} & \textbf{Limitations} \\
    \midrule
    DDPG & Off-policy & Deterministic & Additive OU noise        &
          Sample-efficient & Unstable, $Q$-overestimation \\
    TD3  & Off-policy & Deterministic & Gaussian noise + smoothing &
          Corrects $Q$ bias, stable & Deterministic policy \\
    SAC  & Off-policy & Stochastic    & Max-entropy (implicit)   &
          Robust, natural exploration & $\alpha$ sensitive if poorly initialised \\
    PPO  & On-policy  & Stochastic    & Policy stochasticity     &
          Stable, low variance & Lower sample efficiency \\
    \bottomrule
  \end{tabularx}
  \normalsize
\end{table}

% -----------------------------------------------------------
\section{Reference Controllers (Baselines)}

\subsection{Random Policy}

The random policy draws excitations $u \sim \mathcal{U}([0,1]^9)$ at each step,
without any learning. It establishes the lower performance bound.

\subsection{PD Impedance Controller}

The impedance controller computes a restoring wrench toward the target in
operational space, maps it to joint torques through the transposed site Jacobian,
compensates the bias forces, and distributes the result onto muscles by
minimum-effort load sharing over each muscle's achievable force interval:
\begin{equation}
  \tau^\text{cmd} = J^\top\!\left(K_p (x^\star - x) - K_d \dot{x}\right)
                   + h(q,\dot q),
  \qquad
  \min_{f}\ \sum_i \left(\frac{f_i}{F_{0,i}^M}\right)^{2}
  \ \text{s.t.}\ R^\top f = \tau^\text{cmd},\ f \in [f^{\min}, f^{\max}]
  \label{eq:pd_controller}
\end{equation}
where $h(q,\dot q)$ is \MuJoCo{}'s \texttt{qfrc\_bias} (gravity, Coriolis and
centrifugal terms) and $[f^{\min}, f^{\max}]$ is recovered per muscle by probing
the activation state at $a=0$ and $a=1$.

Gains were selected by grid search, not by hand: a twelve-point coarse grid gave
$K_p = 60$, $K_d = 25$, and a thirty-point refinement gave
$\mathbf{K_p = 20}$, $\mathbf{K_d = 6}$ (\cref{sec:finegains}), which is the
configuration used for every conventional-controller result in this thesis.

Both the bias compensation and the bounded load-sharing step matter, and neither
was present in the first implementation. \Cref{sec:conventional} records what
their absence cost: without gravity compensation the controller reached only
\SI{0.46}{\meter}, worse than a random policy, and an activation mapping that
assumed force proportional to activation produced a zero-width achievable
interval and no motion at all. Both corrections improve the classical side, and
the comparison is reported with them in place.
